\chapter{Networking}
\label{ch:networking}

CNA's \cnans{Microsoft::Xna::Framework::Net} namespace --- five enums and eighteen classes,
mirroring the shape this book's ecosystem chapter already summarized --- is one of the
clearest illustrations in this book of the gap between ``the API surface exists'' and ``the
API surface does something real.'' Real networking is backed by ENet (vendored directly, as
already noted earlier in this book), wrapped by an internal \texttt{ENetBackend}.

\section{NetworkSession: caller-owned, and genuinely capable in ways FNA's stub is not}

\cnaclass{NetworkSession} is explicit about its own ownership model: \texttt{Create()}/%
\texttt{Find()} / \texttt{Join()} / \texttt{JoinInvited()} all return a caller-owned
\cnaclass{NetworkSession*}; \texttt{Dispose()} releases the transport and every gamer the
session owns, but does not delete the session object itself --- the caller must do that
separately, matching the conventional \texttt{IDisposable} split between resource release and
memory ownership. Two properties are worth calling out specifically because they are real in
CNA in a way they are not in FNA at all: \texttt{AllowHostMigration}, where every surviving
peer independently computes the identical new host (the lowest remaining wire ID, requiring
no election round-trip) and reconnects via real LAN discovery; and
\texttt{SimulatedLatency} / \texttt{SimulatedPacketLoss}, which really do delay or
probabilistically drop \texttt{AppData} traffic on the receiving side, scoped specifically to
application data rather than session-management traffic or host-relay hops. In FNA, both of
these properties exist only as stored, inert values with no behavior behind them at all ---
here, they are genuinely functional network-testing tools.

A subtle but load-bearing detail belongs right alongside those two: real XNA's own
\texttt{GamerJoined} is documented to replay itself immediately upon subscription for every
gamer already in the session at that point --- the initial local gamer(s) a
\texttt{Create()} / \texttt{Join()} / \texttt{JoinInvited()} call establishes are queued
internally before the caller could possibly have subscribed to anything, since the session
pointer does not exist until the static factory method actually returns. This port's own
event-handler type has no ``replay on subscribe'' hook to reproduce that automatically, so the
documented, permanent pattern --- not a temporary workaround --- is to call \texttt{Update()}
once, immediately after subscribing, specifically to receive those initial join events.

\begin{lstlisting}
// Worked example: hosting a session and correctly receiving the local
// gamer's own initial GamerJoined event, which is queued at Create()
// time, before any subscription could possibly exist yet.
NetworkSession* session = NetworkSession::Create(
    NetworkSessionType::SystemLink, /* maxLocalGamers */ 1, /* maxGamers */ 8);

session->GamerJoined.Add(
    [](System::Object* sender, const GamerJoinedEventArgs& args)
    {
        // Handle a gamer (local or remote) joining the session.
    });
session->Update();  // required once, right here -- see the note above

// ... in the game loop, every frame:
session->Update();
\end{lstlisting}

\section{NetworkGamer: real per-machine identity where FNA hardcodes it}

\cnaclass{NetworkGamer} / \cnaclass{LocalNetworkGamer}'s \texttt{Id} and \texttt{IsHost} are
real, cross-machine-consistent values in CNA --- a genuine improvement over FNA's own
documented upstream limitation, where \emph{every} gamer's \texttt{Id} is hardcoded to zero
and \emph{every} gamer's \texttt{IsHost} is hardcoded to true, regardless of which machine is
asking. \texttt{RoundtripTime} is wired to the real underlying ENet peer's own native
round-trip-time tracking, not estimated or faked.

\section{Real versus stub, by NetworkSessionType}

This is the single most useful table in this chapter, because it tells you directly whether
a given multiplayer feature in a game you are porting (this book's migration chapter returns
to this) will actually work:

\begin{center}
\begin{tabular}{ll}
\toprule
\textbf{\cnaclass{NetworkSessionType}} & \textbf{Status} \\
\midrule
\texttt{SystemLink} & Real: hosting, joining, LAN discovery, AppData relay, disconnect \\
                     & handling, host migration --- all genuinely networked over ENet \\
\texttt{Local} & Real, loopback-style, no network transport (split-screen/dev testing) \\
\texttt{LocalWithLeaderboards} & \texttt{Local} plus the real, locally-persisted \\
                                & leaderboards/achievements from Chapter~\ref{ch:gamerservices} \\
\texttt{PlayerMatch} / \texttt{Ranked} & Documented stub --- no matchmaking backend exists \\
                                        & to connect to, not an artificially withheld feature \\
\bottomrule
\end{tabular}
\end{center}

Session invites are stubbed for the identical reason \texttt{PlayerMatch} / \texttt{Ranked}
are: there is no online service behind them to implement against, not a deliberately
Xbox-Live-exclusive feature CNA chose not to attempt.

\section{Real LAN discovery, concretely}

\cnaclass{AvailableNetworkSession} / \texttt{AvailableNetworkSessionCollection} are populated by
genuine ENet LAN discovery replies over a well-known UDP port, with broadcast and loopback
fallback and OS-assigned ephemeral ports for the actual game-session transport once a session
is joined. Several accessor methods on \cnaclass{AvailableNetworkSession} --- resolving the
real connect address, port, and session type --- exist purely because CNA's own
\texttt{Find()} genuinely discovers sessions over the network; they have no FNA counterpart
at all, since FNA's own \texttt{Find()} never discovers anything real to begin with.
\cnaclass{QualityOfService} correspondingly has two construction paths: one matching FNA's
own acknowledged-incomplete, no-measurement stub exactly (byte-for-byte, for API parity), and
a second that computes a real round-trip-time measurement from the discovery query/reply
exchange itself --- bandwidth remains unmeasured in both cases, since no established peer
connection exists yet at discovery time to measure it against.

\section{Real bugs found and fixed}

Several genuine, consequential bugs were found and fixed in this namespace's own history,
worth naming because of what each reveals. \cnaclass{NetworkSession}'s own
\texttt{Dispose()} could be called twice and use-after-free as a result --- fixed to be
idempotent. The async \texttt{Begin*} family's callbacks were stored but, for a period, never
actually invoked at all --- fixed to fire exactly once, correctly handling the case where the
callback itself re-entrantly triggers another async call. Every \cnaclass{NetworkGamer}/%
\cnaclass{LocalNetworkGamer} a session ever created used to leak permanently, because nothing
in the session's own lifecycle ever deleted them --- fixed by having \cnaclass{NetworkSession}
take real ownership and free them in bulk on disposal. A local-gamer-ID collision bug could
occur specifically after a remove-then-add sequence, because the ID was derived from the
live collection's current size rather than a monotonically increasing counter --- fixed with a
dedicated counter that never reuses a value. And application data sent before a peer's
handshake had fully completed was silently and permanently dropped, despite being reachable
through the ordinary public send API --- fixed with a small, bounded, per-session queue that
flushes once both sender and target resolve a real wire ID, then further hardened in a later
pass after a follow-up audit found the queue itself was never purged when a gamer left,
disconnected, or a host migration occurred.

One bug remains explicitly open at the time of writing rather than silently dropped: a real,
confirmed leak in \texttt{NetworkSession::BeginCreate}, flagged directly and not yet
resolved.

\section{Two asymmetries, preserved rather than symmetrized}

Two small API details are worth knowing specifically because they look like bugs and are
not: \cnaclass{PacketReader} / \cnaclass{PacketWriter} read and write \cnaclass{Color} values
asymmetrically --- as four floats on read, four bytes on write --- a known, preserved mismatch
matching FNA's own upstream behavior exactly rather than a CNA-introduced inconsistency. And
\cnaclass{NetworkSessionProperties}, a C\texttt{++} translation of a C\# list type whose
single mutable indexer cannot distinguish a read access from a write access the way C\#'s
separate getter/setter accessors can, triggers FNA's own ``auto-append past the end'' growth
behavior on \emph{any} out-of-range access through its mutable overload, not only on an
explicit write --- a documented, accepted structural deviation; a \texttt{const} reference
gives genuinely bounds-checked reads instead.

\begin{lstlisting}
// Worked example: sending a Color over the wire, and the asymmetry above
// caught in the act -- Write(Color) is 4 bytes, but the matching read
// must be ReadColor() (4 floats), not a naive byte-for-byte mirror.
PacketWriter writer;
writer.Write(playerTintColor);      // 4 bytes on the wire
localGamer->SendData(writer, SendDataOptions::Reliable);

// On the receiving end, in the game loop:
PacketReader reader;
NetworkGamer* sender = nullptr;
while (localGamer->ReceiveData(reader, sender) > 0)
{
    Color tint = reader.ReadColor();  // 4 floats -- NOT the inverse of Write(Color) above
}
\end{lstlisting}

\section{Status}

The namespace's own primary hardening plan is fully checked off. Three independent
post-completion audits each found that an earlier ``done'' declaration had overstated
reality --- mostly around Avatar-adjacent rendering concerns (Chapter~\ref{ch:avatar}) rather
than this namespace's own networking logic --- and the specific networking-side gap those
audits did surface (the pre-handshake application-data lifecycle above) is now fixed and
test-confirmed, with a fourth, final audit confirming the networking side as genuinely
settled. Of a broader run of 4{,}935 tests touching this area, 4{,}884 pass; every one of the
51 remaining failures has been individually traced to a missing test fixture, a software
rendering limitation, or parallel-test port contention on the well-known discovery port ---
none traced back to this namespace's own logic.
